% =====================================================================
%  IEEE-STYLE CONFERENCE REPORT TEMPLATE  (self-contained)
%  MLE-AI-201 Machine Learning -- Final Project
%
%  This template MIMICS the IEEE two-column conference look using only
%  standard LaTeX packages. It does NOT depend on IEEEtran.cls, so it
%  compiles anywhere (tectonic, pdflatex, Overleaf) with nothing extra.
%
%  Build:   tectonic report.tex      (recommended)
%     or:   pdflatex report.tex ; pdflatex report.tex
% =====================================================================
\documentclass[10pt,twocolumn,a4paper]{article}

% ---- Geometry: IEEE-like margins / two narrow columns ----------------
\usepackage[a4paper,top=0.75in,bottom=1.0in,left=0.68in,right=0.68in,columnsep=0.25in]{geometry}

% ---- Times-like font (standard, ships with every TeX distro) ---------
\usepackage{mathptmx}          % Times text + math
\usepackage[T1]{fontenc}
\usepackage[utf8]{inputenc}    % allows accented / Azerbaijani chars

% ---- Common, dependency-free packages --------------------------------
\usepackage{graphicx}
\usepackage{amsmath,amssymb}
\usepackage{booktabs}
\usepackage{array}
\usepackage{titlesec}
\usepackage{caption}
\usepackage{cite}
\usepackage[hidelinks]{hyperref}
\usepackage{url}

% ---- IEEE-style section headings -------------------------------------
\renewcommand{\thesection}{\Roman{section}}
\renewcommand{\thesubsection}{\Alph{subsection}}
\titleformat{\section}[block]{\centering\normalsize\scshape}{\thesection.}{0.5em}{}
\titleformat{\subsection}[block]{\normalsize\itshape}{\thesubsection.}{0.5em}{}
\titlespacing*{\section}{0pt}{1.4ex plus 1ex minus .2ex}{1ex}
\titlespacing*{\subsection}{0pt}{1.0ex plus 1ex minus .2ex}{0.6ex}

% ---- Abstract / Index-Terms in the IEEE boldface-dash style ----------
\newcommand{\abstractkw}[1]{\noindent\textbf{\textit{#1}}}

\captionsetup{font=small,labelsep=period}
\setlength{\parindent}{1em}

% =====================================================================
\begin{document}

% ------------------------- TITLE BLOCK --------------------------------
\twocolumn[
  \begin{@twocolumnfalse}
  \begin{center}
    {\LARGE\bfseries Real-Estate Price Prediction:\\Decision Trees and Support Vector Machines\par}
    \vspace{1.2em}
    % Authors: Given name Surname, comma-separated.
    {\large
      Leyla Aliyeva, Agil Asgarli, Mahammad Gurbanzade, Dilara Rzazade\par}
    \vspace{0.5em}
    {\itshape MLE-AI-201 Machine Learning, AI Academy --- Cohort II\par}
    \vspace{1.4em}
  \end{center}
  \begin{quote}
    \small
    \abstractkw{Abstract---}% IEEE uses an em-dash after "Abstract"
    We predict real-estate asking prices and classify listings into standard and premium tiers using bina.az data. Our pipeline cleans duplicate listings, handles missing values, encodes categorical variables, and prevents target leakage. We implement decision-tree models and a linear SVM trained with Pegasos, then compare them with scikit-learn baselines. After cleaning, 64,452 listings and 173 features remain. The premium threshold is the training-set median of 210,000 AZN. Our decision-tree classifier achieves 0.9180 test accuracy and 0.9588 ROC-AUC; the regression tree achieves an original-scale R-squared of 0.5144. The linear SVM reaches 0.8492 test accuracy and 0.9233 ROC-AUC, comparable to scikit-learn's linear SVMs but about seven points below the tree. These results demonstrate the usefulness of tree-based models for property listings, while the use of asking prices and a random split limits conclusions about future market prices. Code: \href{https://github.com/Dilara898/ML-Final_Project}{\texttt{github.com/Dilara898/ML-Final\_Project}}.

    \vspace{0.6em}
    \abstractkw{Index Terms---}decision tree, support-vector machine,
    Pegasos, regression, classification, real-estate price prediction.
  \end{quote}
  \vspace{1.0em}
  \end{@twocolumnfalse}
]

% =====================================================================
\section{Introduction}
Real-estate price estimation can help compare listings and understand how property attributes relate to asking prices. Prices depend on location, area, room count, and building characteristics, making this a useful prediction problem. This project addresses regression on log-transformed prices and classification into standard and premium tiers. Our contributions are a reproducible preprocessing pipeline, from-scratch NumPy decision-tree and Pegasos SVM implementations, comparison with scikit-learn baselines, and held-out test evaluation. Training-only preprocessing and explicit leakage prevention are used to make the evaluation more reliable.

\section{Related Work / Background}
Decision trees are supervised learning models that recursively partition
the feature space. Classification trees commonly select splits by reducing
impurity, measured using criteria such as Gini impurity or entropy.
Regression trees instead seek splits that reduce squared prediction error.
Their main advantage is the ability to represent nonlinear relationships
and feature interactions without requiring a linear functional form.
However, overly deep trees may overfit, so model complexity must be controlled.

Support Vector Machines (SVMs) seek a decision boundary with a large margin.
The soft-margin formulation balances margin size against hinge-loss
classification errors through a regularisation parameter. Pegasos optimises
the regularised hinge-loss objective using stochastic sub-gradient updates,
making linear SVM training practical on large datasets~\cite{pegasos}.
Random Fourier Features (RFF) approximate an RBF kernel through a randomized
feature mapping, allowing a linear learner to represent nonlinear boundaries
without constructing a full kernel matrix.

\section{Data and Preprocessing}
The dataset was collected from bina.az real-estate listings. After cleaning,
64,452 rows and 173 model input features remained. The data were divided into
45,116 training examples, 9,667 validation examples, and 9,669 test examples.
The training set contains 49.37\% premium listings. The premium threshold is
210,000 AZN, defined using the training-set median.

Two tasks were considered. Regression predicts the asking price (modelled
as $\log(\text{price})$ because prices are heavy-tailed), while classification
predicts whether a listing belongs to the premium tier.

The raw dump contained 100{,}775 rows and 51 columns. Repeated observations of
the same listing were identified by listing URL (estate IDs were not unique per
property); we kept the most recent scrape of each URL, which removed 36{,}321
rows and left 64{,}454. Two listings whose sale-price target was ambiguous were
excluded after manual review, giving 64{,}452 rows. Both currency columns
contained AZN throughout, so no conversion was needed. Area values recorded in
\emph{sot} were converted to square metres, and floor strings such as ``7\,/\,9''
were split into floor number and total floors. Missing renovation and mortgage
information was encoded as ``unknown'' rather than assumed to mean ``no''.
Remaining numerical gaps were filled with training-set medians within each
property category, with missingness indicators for land area, floor, total
floors, rooms and coordinates; such filled values are modelling placeholders,
not verified measurements. The building-type column was dropped because about
99.85\% of its values were missing.

The columns \texttt{unit\_price} and \texttt{total\_price} restate the target
and were removed to prevent target leakage; identifiers, URLs, owner details
and free text were not used. Model inputs come from an explicit list of 18
columns, which expand to 173 features after categorical encoding. Training rows
(about 70\%) were drawn at random with seed 42, the tier threshold was computed
from the training median, and the remaining rows were split into validation and
test sets stratified by tier. Scaling statistics come from the training split
only. Upper limits for \texttt{area\_m2} and \texttt{land\_area\_m2} are the 99th training percentile within each property category (learned only for categories with at least 100 observed training values); no rows were removed and prices were not capped. One listing whose coordinates conflicted with its stated location had them set to missing. No listing URL appears in more than one split. Different URLs may still describe the same physical property, so a
random split can place near-duplicates in different splits.

Exploratory analysis of the cleaned training data (before capping and imputation)
shows a long right tail in prices, so we model $\log(\text{price})$; the plot
alone does not establish a log-normal distribution. Area is positively
associated with price, with wide variation within each property category, and
some unusually low prices and large areas remain. Two- and three-room
properties dominate the residential listings, and median price generally rises
with room count, although area, location and renovation also differ between
properties, so rooms alone do not explain the trend. The data are heavily
concentrated in Baku listings, so the results mainly describe that market. The
median threshold gives nearly balanced tiers: 22,840 standard and 22,276
premium listings in training (49.37\% premium) and 49.45\% premium in both the
validation and test splits. Prices are advertised asking prices, not verified
transaction prices.

\begin{table}[t]
  \centering
  \caption{Dataset summary after cleaning.}
  \label{tab:data}
  \begin{tabular}{@{}lr@{}}
    \toprule
    Property & Value \\
    \midrule
    Rows (after cleaning)       & 64,452 \\
    Features used               & 173 \\
    Train / Val / Test split    & 45,116 / 9,667 / 9,669 \\
    Premium share (train)       & 49.37\% \\
    \bottomrule
  \end{tabular}
\end{table}

\section{Method}
\subsection{Decision tree (from scratch)}
The tree recursively selects feature splits that reduce node impurity for classification or squared prediction error for regression. Each example follows the learned rules to a leaf, which returns a class prediction or numerical estimate. Tree depth and leaf-size constraints control complexity. Models are evaluated on the held-out split and compared with scikit-learn decision-tree baselines.

\subsection{Support vector machine (from scratch)}
We implement a linear soft-margin SVM in NumPy and train it with Pegasos,
a stochastic sub-gradient method on the regularised hinge-loss primal
(no quadratic-programming solver). Labels are mapped to $\{-1,+1\}$
(premium $=+1$). At step $t$ we draw one example $i$, set
$\eta_t = 1/(\lambda t)$, and update
\begin{equation}
  \mathbf{w} \leftarrow (1-\eta_t\lambda)\,\mathbf{w}
  + \eta_t\, y_i \mathbf{x}_i \,\mathbf{1}\!\left[y_i\,\mathbf{w}^{\top}\mathbf{x}_i < 1\right].
\end{equation}
The bias is folded into $\mathbf{w}$ by appending a constant 1 to every
feature vector, so the bias is also regularised (a difference from the
standard formulation and from scikit-learn). Features are standardised
with training-set statistics only. The margin width is $2/\lVert\mathbf{w}\rVert$
(computed on the weights without the bias), and support vectors are the
training points with $y_i(\mathbf{w}^{\top}\mathbf{x}_i+b)\le 1$. The
objective $J$ is logged during training. A constant learning-rate option
is also implemented and compared below; all other reported results use the Pegasos schedule.

For non-linear boundaries we use random Fourier features (RFF) rather than
kernelised Pegasos: $z(\mathbf{x})=\sqrt{2/D}\cos(W^{\top}\mathbf{x}+\mathbf{b})$
with $W_{jk}\sim\mathcal{N}(0,2\gamma)$ and $b_k\sim U[0,2\pi]$, fitted on
the training split only. This approximates an RBF kernel, reuses the linear
solver unchanged, and avoids an $n\times n$ kernel matrix.

\subsection{SVM tuning and results}
SVM hyperparameters were chosen on the validation split only; the test set
was evaluated once, after the final configuration was fixed. We swept
$\lambda\in\{10^{-5},\dots,1\}$ (50{,}000 iterations), then, for RFF,
$\gamma\in\{0.0005,\dots,0.03\}$ with $\lambda\in\{10^{-5},10^{-4},10^{-3}\}$
and $D=300$. A follow-up checked longer training (50k, 200k, 500k iterations)
and $D\in\{300,1000\}$. All runs use seed 42. Baselines are scikit-learn's
SGD classifier (hinge loss) and LinearSVC on standardised features.

Table~\ref{tab:svm_lambda} and Fig.~\ref{fig:svm_conv} show the effect of $\lambda$: as $\lambda$ grows,
$\lVert\mathbf{w}\rVert$ shrinks and the margin widens, while validation
accuracy peaks at $\lambda=10^{-2}$ (with 50k iterations).

\begin{table}[t]
  \centering
  \caption{Linear Pegasos, $\lambda$ sweep (50k iterations $\approx$ 1.1 epochs).}
  \label{tab:svm_lambda}
  \begin{tabular}{@{}rrrrrr@{}}
    \toprule
    $\lambda$ & $\lVert w\rVert$ & margin & SV share & train & val \\
    \midrule
    $10^{-5}$ & 788.8 & 0.003 & 0.208 & 0.794 & 0.795 \\
    $10^{-4}$ & 79.9  & 0.025 & 0.222 & 0.801 & 0.801 \\
    $10^{-3}$ & 9.83  & 0.203 & 0.291 & 0.813 & 0.818 \\
    $10^{-2}$ & 2.37  & 0.845 & 0.408 & 0.838 & 0.844 \\
    $10^{-1}$ & 1.12  & 1.781 & 0.542 & 0.826 & 0.827 \\
    $1$       & 0.54  & 3.713 & 0.855 & 0.792 & 0.790 \\
    \bottomrule
  \end{tabular}
\end{table}

\begin{figure}[t]
  \centering
  \includegraphics[width=\columnwidth]{figures/svm_lambda_convergence.png}
  \caption{Pegasos objective $J$ versus iteration for each $\lambda$ (log scale).
  Small $\lambda$ ($\le 10^{-3}$) is still decreasing after 50k iterations, whereas $\lambda \ge 10^{-2}$ is essentially flat.}
  \label{fig:svm_conv}
\end{figure}

With longer training (500k iterations, about 11 epochs), $\lambda=3\times10^{-3}$
gave the best validation accuracy (0.846) and was fixed as the final
configuration. On the test set the SVM reaches accuracy 0.849, F1 0.850 and
AUC 0.923; in 5-fold CV its F1 is $0.843\pm0.002$, versus $0.845\pm0.006$ for
SGD-hinge and $0.839\pm0.003$ for LinearSVC, i.e.\ within one standard
deviation of both baselines. The final model has $\lVert w\rVert=2.52$, margin
width 0.79, and 17{,}288 support vectors (38.3\% of training data). Fitting
takes about 6.4\,s (SGD-hinge 1.1\,s, LinearSVC 3.6\,s) and prediction about 1\,ms
(Table~\ref{tab:timing}).

\paragraph{Learning-rate schedule.} At the final $\lambda=3\times10^{-3}$ and 500k iterations, the Pegasos schedule $\eta_t=1/(\lambda t)$ reached the lowest objective ($J=0.390$) and the best validation accuracy (0.846). A small constant step ($\eta_0=0.001$) was almost as good ($J=0.395$, 0.844), but larger constant steps degraded the solution: $\eta_0=0.01$ gave $\lVert w\rVert=4.24$, margin 0.47 and validation accuracy 0.825, and $\eta_0=0.1$ gave $\lVert w\rVert=15.1$, margin 0.13, $J=1.87$ and accuracy 0.782 (Table~\ref{tab:svm_lr}). We therefore keep the Pegasos schedule.

\begin{table}[t]
  \centering
  \footnotesize
  \caption{Learning-rate schedules at $\lambda=3\times10^{-3}$, 500k iterations.}
  \label{tab:svm_lr}
  \begin{tabular}{@{}lrrrrrr@{}}
    \toprule
    schedule & $\lVert w\rVert$ & margin & SV share & train & val & final $J$ \\
    \midrule
    Pegasos & 2.52 & 0.793 & 0.383 & 0.845 & 0.846 & 0.390 \\
    constant 0.001 & 2.54 & 0.788 & 0.380 & 0.843 & 0.844 & 0.395 \\
    constant 0.01 & 4.24 & 0.472 & 0.341 & 0.828 & 0.825 & 0.534 \\
    constant 0.1 & 15.1 & 0.132 & 0.264 & 0.787 & 0.782 & 1.867 \\
    \bottomrule
  \end{tabular}
\end{table}

RFF did not improve on the linear SVM: the best validation accuracy was 0.830
with $D=300$ and 0.843 with $D=1000$, versus 0.846 for the linear model.
Increasing the number of iterations did not help at $D=300$, whereas more
components did, which suggests the feature budget, not training length, was
limiting; we did not test this further. The decision tree outperforms the
SVM by about seven accuracy points; a plausible reason (not verified) is that
a linear model cannot capture interactions between, for example, location
and area. Differences of about 0.004 in validation accuracy are within
sampling noise for 9{,}667 validation examples, so close configurations
cannot be ranked reliably. Remaining differences to scikit-learn come from
the stochastic solver, the regularised bias, and the $C=1/(n\lambda)$
correspondence being exact only at the optimum. Sweeps used a single seed.


\section{Experimental Setup and Results}
\subsection{Protocol and metrics}
Task~A is evaluated with RMSE, MAE and $R^2$, both on $\log(\text{price})$
(the training target) and on the original AZN scale. Task~B (premium $=1$) is
evaluated with accuracy, precision, recall, F1 and ROC-AUC. All metrics are
implemented in NumPy and unit-tested against scikit-learn. Final numbers use
the test split (9{,}669 listings), on which models were fitted on the training
split only. Uncertainty is the mean $\pm$ sample standard deviation over five
folds of the training and validation rows; the folds are identical for our
models and the baselines, stratified by tier, and the tier threshold and
scaler are recomputed inside every fold.

Our classification tree uses entropy with depth 20 and
\texttt{min\_samples\_leaf}\,=\,5; our regression tree uses squared error
with depth 20 and \texttt{min\_samples\_leaf}\,=\,20; both use
\texttt{min\_samples\_split}\,=\,10. The baselines are scikit-learn decision
trees with the same settings, an SGD classifier (hinge loss, $\alpha=\lambda$)
and LinearSVC ($C=1/(n\lambda)$), the last two with a \texttt{StandardScaler}
inside a \texttt{Pipeline}. An RBF \texttt{SVC} was not run because its cost
grows quickly with the 45{,}116 training rows. The whole benchmark is produced
by \texttt{python -m src.run\_all}; all seeds are fixed to 42, and the
\texttt{--check-seeds} option confirmed that two complete runs produce
identical test metrics.

\subsection{Comparison with scikit-learn}
Table~\ref{tab:testA} and Table~\ref{tab:testB} give the test results. Our
trees are practically indistinguishable from scikit-learn's: classification
accuracy differs by 0.0006, ROC-AUC by less than 0.0001, and the log-scale RMSE by
0.0006. We did not isolate the cause of these small gaps; tie-breaking among
equally good splits and the exact threshold convention are plausible sources.
The regression tree explains 77.4\% of the variance of $\log(\text{price})$ but
only 51.4\% on the AZN scale, where a few very expensive listings dominate the
squared error (the most expensive test listing costs 18.5 million AZN). For
the same reason the price-scale RMSE is unstable across cross-validation folds
(about $1.68\times10^{6}\pm2.28\times10^{6}$ AZN), so we compare regression
models in cross-validation on the log scale: RMSE $0.390\pm0.011$ and $R^2$
$0.767\pm0.013$ for both our tree and scikit-learn's.

\begin{table}[t]
  \centering
  \footnotesize
  \caption{Task A on the test set. Price-scale errors are in AZN.}
  \label{tab:testA}
  \begin{tabular}{@{}lccccc@{}}
    \toprule
    Model & RMSE & MAE & $R^2$ & RMSE$_{\log}$ & $R^2_{\log}$ \\
    \midrule
    Our tree     & 444{,}189 & 89{,}583 & 0.5144 & 0.3835 & 0.7740 \\
    sklearn tree & 444{,}104 & 89{,}553 & 0.5146 & 0.3829 & 0.7746 \\
    \bottomrule
  \end{tabular}
\end{table}

\begin{table}[t]
  \centering
  \footnotesize
  \caption{Task B on the test set (premium $=1$).}
  \label{tab:testB}
  \begin{tabular}{@{}lccccc@{}}
    \toprule
    Model & Acc. & Prec. & Rec. & F1 & AUC \\
    \midrule
    Our tree     & 0.9180 & 0.9168 & 0.9174 & 0.9171 & 0.9588 \\
    sklearn tree & 0.9174 & 0.9158 & 0.9172 & 0.9165 & 0.9588 \\
    Our SVM      & 0.8492 & 0.8371 & 0.8630 & 0.8498 & 0.9233 \\
    SGD-hinge    & 0.8528 & 0.8392 & 0.8689 & 0.8538 & 0.9287 \\
    LinearSVC    & 0.8438 & 0.8260 & 0.8668 & 0.8459 & 0.9269 \\
    \bottomrule
  \end{tabular}
\end{table}

\begin{table}[t]
  \centering
  \footnotesize
  \caption{Task B, 5-fold cross-validation on train+val (mean $\pm$ std).}
  \label{tab:cv}
  \begin{tabular}{@{}lcc@{}}
    \toprule
    Model & F1 & ROC-AUC \\
    \midrule
    Our tree     & $0.915\pm0.004$ & $0.955\pm0.003$ \\
    sklearn tree & $0.915\pm0.004$ & $0.954\pm0.003$ \\
    Our SVM      & $0.843\pm0.002$ & $0.923\pm0.004$ \\
    SGD-hinge    & $0.845\pm0.006$ & $0.926\pm0.004$ \\
    LinearSVC    & $0.839\pm0.003$ & $0.925\pm0.003$ \\
    \bottomrule
  \end{tabular}
\end{table}

The linear SVMs, ours included, are about seven accuracy points behind the
trees, and the cross-validated differences between the three linear models are
within about one standard deviation. Our implementations are slower than
scikit-learn's (Table~\ref{tab:timing}): our trees take roughly twenty to thirty
times longer to fit, plausibly because scikit-learn's trees are compiled code
whereas ours run in NumPy and Python; we did not profile this. The timings come
from a single run, so they indicate magnitude only.

\begin{table}[t]
  \centering
  \footnotesize
  \caption{Fit and prediction time on 45{,}116 training rows (single run).}
  \label{tab:timing}
  \begin{tabular}{@{}lcc@{}}
    \toprule
    Model & Fit (s) & Predict (ms) \\
    \midrule
    Our tree (regression)      & 10.96 & 36.2 \\
    sklearn tree (regression)  & 0.50 & 4.9 \\
    Our tree (classification)  & 15.73 & 36.2 \\
    sklearn tree (classif.)    & 0.49 & 5.9 \\
    Our SVM (Pegasos)          & 6.41 & 1.1 \\
    SGD-hinge                  & 1.09 & 10.9 \\
    LinearSVC                  & 3.62 & 18.8 \\
    \bottomrule
  \end{tabular}
\end{table}

\subsection{Error analysis}
Fig.~\ref{fig:error_decile} groups test listings by true-price decile. The
median absolute percentage error of the regression tree is between 7.8\% and
11.5\% for deciles D2--D8, rises to 12.0\% in D9 and to 25.1\% in D10, and is
22.6\% in D1. The most expensive decile alone accounts for 63\% of the total
absolute price error, with a mean absolute error of about 553{,}000\,AZN
against 21{,}000--79{,}000\,AZN in the other deciles, so the model is least
reliable for luxury listings. The cheapest decile starts at 127\,AZN, which
suggests that some entries are data errors; we did not check this.

For classification, the tree makes 398 false positives and 395 false negatives
on the test set (Fig.~\ref{fig:cm_tree}), and the SVM makes 803 and 655
(Fig.~\ref{fig:cm_svm}). The two models disagree on 15.4\% of test listings;
the tree is right and the SVM wrong on 1{,}075 of them, the reverse on 410.

The tree importances (Fig.~\ref{fig:importance_clf}) are dominated by floor
area (0.48 for the classifier, 0.56 for the regressor), followed by longitude
and latitude and, further down, the land-plot category, land area and the
new-building category. Impurity-based
importances tend to favour continuous features with many split points, so they
indicate how the fitted tree uses a feature rather than a causal effect. The
SVM weights (Fig.~\ref{fig:svm_weights}) are largest for the number of rooms
(1.69) and floor area (0.94), positive for the new-building category and
negative for the old-building category and for several peripheral districts.
Rooms and area are strongly related, so their individual weights should not be
interpreted separately.

\section{Visual Results}
The following figures provide complementary views of classification
performance, feature relevance, and regression error.

\begin{figure}[t]
\centering
\includegraphics[width=\columnwidth]{figures/roc_curves.pdf}
\caption{ROC curves comparing classification models. A larger AUC indicates
better separation between the two classes across decision thresholds.}
\label{fig:roc}
\end{figure}

\begin{figure}[t]
\centering
\includegraphics[width=\columnwidth]{figures/cm_tree.pdf}
\caption{Confusion matrix for the decision-tree classifier. The diagonal
contains correct predictions; off-diagonal entries represent errors.}
\label{fig:cm_tree}
\end{figure}

\begin{figure}[t]
\centering
\includegraphics[width=\columnwidth]{figures/importance_tree_clf.pdf}
\caption{Feature importance for the decision-tree classifier. Importance
scores describe the contribution of features to the fitted tree's splits.}
\label{fig:importance_clf}
\end{figure}

\begin{figure}[t]
\centering
\includegraphics[width=\columnwidth]{figures/pred_vs_actual_tree.pdf}
\caption{Predicted versus actual prices for the regression tree. Points
closer to the ideal diagonal correspond to more accurate predictions.}
\label{fig:pred_actual}
\end{figure}

\begin{figure}[t]
\centering
\includegraphics[width=\columnwidth]{figures/residuals_tree.pdf}
\caption{Regression residuals. Patterns in residuals can indicate systematic
errors that a model does not capture.}
\label{fig:residuals}
\end{figure}

\begin{figure}[t]
\centering
\includegraphics[width=\columnwidth]{figures/error_by_price_decile.pdf}
\caption{Prediction error across actual-price deciles, showing how errors
vary between less expensive and more expensive properties.}
\label{fig:error_decile}
\end{figure}

\begin{figure}[t]
\centering
\includegraphics[width=\columnwidth]{figures/cm_svm.pdf}
\caption{Confusion matrix for the SVM classifier.}
\label{fig:cm_svm}
\end{figure}

\begin{figure}[t]
\centering
\includegraphics[width=\columnwidth]{figures/svm_weights.pdf}
\caption{Learned SVM weights. Their magnitude indicates the model's
linear coefficients, subject to the feature scaling used during training.}
\label{fig:svm_weights}
\end{figure}


\clearpage
\section{Conclusion}
The decision-tree classifier achieved 0.9180 test accuracy, an F1 score of 0.9171, and ROC-AUC of 0.9588. The from-scratch linear SVM reached 0.8492 test accuracy, comparable to scikit-learn's linear SVMs but about seven points below the tree; random Fourier features did not close the gap. The regression tree achieved an R-squared of 0.5144 on the original price scale, leaving substantial price variation unexplained. Trees can capture nonlinear relationships, while the linear SVM provides a useful contrasting baseline. Future work could add richer location features and test on a time-based split. Predictions concern asking prices, not confirmed sale prices.

\section*{AI-Assistance Disclosure}
An AI assistant (Claude, Anthropic) was used in two parts of the project. For
the SVM, it explained Pegasos and random Fourier features and drafted most of the code in
\texttt{src/svm.py}, the sweep scripts, the SVM tests and the SVM sections of this report. For evaluation and
integration, it drafted \texttt{src/evaluate.py} (metrics, ROC curve, k-fold
splitting), the figure helpers, \texttt{src/run\_all.py}, the associated tests,
and the experimental-setup, results and error-analysis text of this report.
Team members ran and checked this output: the metrics are tested against
scikit-learn, and the pipeline was run end to end on the real dataset with a
fixed seed. The decision tree and the SVM are our own NumPy implementations;
scikit-learn is used only for baselines and as a reference in tests.

\begin{thebibliography}{9}
\bibitem{pegasos}
S.~Shalev-Shwartz, Y.~Singer, N.~Srebro, and A.~Cotter,
``Pegasos: primal estimated sub-gradient solver for SVM,''
\emph{Mathematical Programming}, vol.~127, no.~1, pp.~3--30, 2011.
\end{thebibliography}

\end{document}


